\section{Rodando o Primeiro Container e Comandos de Sobrevivência}

\begin{frame}
  \frametitle{Rodando o Primeiro Container}
  Após a instalação, execute:

  \vspace{0.5em}
  \texttt{docker run -d -p 8080:80 docker/welcome-to-docker}

  \vspace{0.5em}
  \begin{itemize}
    \item Frontend acessível em \texttt{http://localhost:8080}
    \item Container aparece como rodando
    \item Image passa a constar nas images locais
  \end{itemize}
\end{frame}

\begin{frame}
  \frametitle{Executando Containers}
  \texttt{docker run -d -p 8080:80 docker/welcome-to-docker}

  \vspace{0.5em}
  \begin{itemize}
    \item \texttt{-d} — executa em modo \textit{detached} (segundo plano)
    \item \texttt{-p 8080:80} — mapeia a porta 8080 do \textit{host} para a porta 80 do container
  \end{itemize}
\end{frame}

\begin{frame}
  \frametitle{Listando e Parando Containers}
  \texttt{docker ps} \hspace{1em} — lista containers em execução

  \vspace{0.3em}
  \texttt{docker ps -a} \hspace{0.5em} — lista todos, incluindo os parados

  \vspace{0.8em}
  \texttt{docker stop <ID>} — parar o container

  \vspace{0.5em}
  \begin{block}{Dica}
    Não é necessário digitar o ID completo — basta os caracteres suficientes
    para torná-lo único.
  \end{block}
\end{frame}

\begin{frame}
  \frametitle{Removendo Containers}
  \texttt{docker rm <ID>} \hspace{1.5em} — remove um container parado

  \vspace{0.5em}
  \texttt{docker rm -f <ID>} \hspace{0.5em} — para e remove (força)

  \vspace{0.8em}
  A flag \texttt{-f} força a remoção mesmo que o container ainda esteja em execução — ele é parado automaticamente antes de ser removido.
\end{frame}

\begin{frame}
  \frametitle{Visualizando Logs}
  \texttt{docker logs <ID>} \hspace{2.2em} — exibe todos os logs

  \vspace{0.3em}
  \texttt{docker logs -f <ID>} \hspace{1.2em} — segue logs em tempo real

  \vspace{0.3em}
  \texttt{docker logs --tail 50 <ID>} — últimas 50 linhas

  \vspace{0.8em}
  \begin{itemize}
    \item Exibe \texttt{stdout} e \texttt{stderr} do container
    \item A flag \texttt{-f} funciona como \texttt{tail -f} no Linux
    \item Útil para monitorar o comportamento da aplicação em tempo real
  \end{itemize}
\end{frame}

\begin{frame}
  \frametitle{Executando Comandos Dentro de um Container}
  \texttt{docker exec -it <ID> /bin/bash} — abre shell interativo

  \vspace{0.3em}
  \texttt{docker exec <ID> ls /app} \hspace{1em} — executa comando único

  \vspace{0.3em}
  \texttt{docker exec -it <ID> sh} \hspace{1.5em} — alternativa para Alpine

  \vspace{0.8em}
  \begin{itemize}
    \item \texttt{-i} — mantém o stdin aberto (interativo)
    \item \texttt{-t} — aloca um terminal (TTY)
    \item Juntas, \texttt{-it} permitem interagir como um terminal remoto
  \end{itemize}

  \vspace{0.5em}
  \begin{block}{Uso}
    Depuração, inspeção de arquivos, verificação de configurações ou
    comandos administrativos dentro do container.
  \end{block}
\end{frame}

\begin{frame}
  \frametitle{Resumo dos Comandos}
  \begin{center}
  \small
  \begin{tabular}{|l|l|}
  \hline
  \textbf{Comando} & \textbf{Função} \\
  \hline
  \texttt{docker run} & Cria e executa um container \\
  \texttt{docker ps / ps -a} & Lista containers \\
  \texttt{docker stop <ID>} & Para um container \\
  \texttt{docker rm <ID>} & Remove um container parado \\
  \texttt{docker rm -f <ID>} & Para e remove \\
  \texttt{docker logs -f <ID>} & Segue logs em tempo real \\
  \texttt{docker exec -it <ID> sh} & Abre shell no container \\
  \hline
  \end{tabular}
  \end{center}
\end{frame}
